\chapter{总结与展望}

\section{全文总结}

本文围绕深度学习免疫多肽结合预测模型的优化与部署展开，以 FennOmix.MHC 为基础，针对其在大规模人类蛋白质组应用中暴露出的高维表示、缓存膨胀、I/O 负担重和部署代价高等问题，从模型轻量化和系统工程优化两个方向进行了探索。

在模型优化方面，我们首先复现了 FennOmix.MHC，随后依次尝试了输出维度压缩、降维结构对比和模型剪枝。实验表明，适度压缩嵌入维度能够在守住核心预测能力的同时，大幅削减缓存规模；对比线性与非线性降维方案后，线性结构在性能、简洁性和下游任务表现之间取得了更均衡的效果；剪枝实验则揭示了模型内部存在冗余，经合理裁剪后模型表现依然稳健。

在系统优化方面，我们针对大规模肽段缓存的存储与访问难题，设计了缓存压缩、索引定位和分阶段加载机制。通过将离线预计算、压缩缓存、哈希索引和在线查询整合为一条完整的优化链路，有效降低了重复编码开销，缓解了缓存存储压力，提升了模型在大规模候选肽段筛选场景中的可用性。

总体来看，我们并没有孤立地追求单一指标的提升，而是紧扣 FennOmix.MHC 的实际部署需求，对模型表示压缩、结构优化、缓存管理和索引检索进行了统一设计。结果表明，面向大规模候选库的免疫多肽结合预测模型，在关注预测精度的同时，必须一并考量存储成本、访问效率和运行环境等工程因素。本文的工作验证了这一优化思路的可行性，并为后续系统的迭代改进提供了参考。

\section{本文工作不足}

尽管围绕模型轻量化与部署优化开展了较系统的研究，当前工作仍存在若干局限。实验设计更侧重验证整体技术路线的可行性，对部分参数组合的探索尚欠充分：剪枝部分已证实基于重要性分析的策略可行，但不同剪枝比例、层级组合以及剪枝后推理速度变化的系统比较仍需细化；缓存压缩方面，虽从误差角度验证了方法的有效性，而不同量化参数对检索结果及系统延迟的影响仍有待深入剖析。此外，系统虽已串接为原型链路，但工程完备性尚有提升空间。模型、缓存、索引与在线查询之间的多模块协同对数据格式、版本管理和接口一致性提出了较高要求，在大规模离线构建或多进程运行场景下，还需补充任务调度、断点续算、异常处理和日志记录等功能。最后，部署分析主要依托单机环境与原型验证，对不同硬件条件、查询负载及高并发场景的性能表现仍缺少更充分的测试，因此当前成果更适合定位为方法验证与系统原型基础，距离成熟稳定的实际应用系统仍有完善余地。

\section{后续研究展望}

针对上述不足，未来工作可从以下几个方向推进。其一，细化模型轻量化策略：在降维方面，继续探索更契合免疫多肽结合预测特性的压缩结构和低位宽表示；在剪枝方面，系统比较不同剪枝比例、位置与组合对模型性能、参数量及推理速度的影响，以形成更为可控的模型压缩方案。其二，完善缓存压缩与索引系统的工程实现：在现有原型上补充缓存一致性校验、异常处理、日志管理与自动化构建流程，并测试不同查询规模下全量加载与延迟加载两种模式的实际表现，使部署分析更为完整。其三，拓展技术路线的适用边界：本文凝练的“模型轻量化 + 缓存压缩 + 索引检索”思路并不局限于 FennOmix.MHC，也可为其他面向大规模候选库筛选的生物信息学深度学习模型提供借鉴。随着蛋白质语言模型、对比学习与向量检索技术在生命科学中的不断渗透，如何在维持表示能力的同时压低系统部署成本，将成为值得持续探究的议题。

\section{本章小结}

本章归纳了全文的研究工作，并剖析了当前存在的不足与后续优化方向。概括而言，本研究瞄准 FennOmix.MHC 在大规模人类蛋白质组场景中的部署瓶颈，完成了模型轻量化、缓存压缩与索引检索等环节的探索，验证了模型优化与系统部署协同设计的可行性。未来仍需在参数精细化、工程鲁棒性提升及多场景部署测试等方面持续完善。